\section{问题分析}
\subsection{总体建模思路}
四个问题并不是彼此独立的计算任务，而是从理想二维剖面逐步过渡到真实三维地形的同一套测线规划问题。本文首先确定声束扇面与海底剖面的几何关系，以局部水深和横航向坡度作为覆盖宽度的最小状态变量；随后把海底坡度分解到航向和横航向，使同一覆盖公式能够处理任意测线方向；在规则矩形海域中，利用覆盖边界和重叠率下限构造可直接计算的递推关系；在真实海域中，则将局部覆盖公式嵌入地形插值、候选路线生成、曲率平滑和独立空间评估。这样处理的好处是，后一个问题只增加必要的新变量和新约束，不重新定义前面已经建立的物理量。

图~\ref{fig:flow} 给出本文的完整求解流程。流程中“路线生成”和“指标评估”被刻意分离：路线生成器只依据地形和重叠规则给出候选几何，独立评估器重新计算覆盖、长度、边界和曲率指标，从而避免用生成算法自身的内部得分直接证明方案有效。

\begin{figure}[H]
\centering
\begin{tikzpicture}[
node distance=7mm and 8mm,
box/.style={draw,rounded corners,align=center,minimum width=31mm,minimum height=8mm,font=\small},
arr/.style={-{Latex[length=2mm]},thick}
]
\node[box] (input) {题面与附件\\参数、网格、约束};
\node[box,right=of input] (p1) {二维斜坡\\覆盖几何};
\node[box,right=of p1] (p2) {方向分解\\等效坡度};
\node[box,below=of p1] (p3) {显式递推\\规则海域布线};
\node[box,right=of p3] (p4) {地形插值与\\有限曲率路线};
\node[box,left=of p3] (check) {独立复算\\边界与敏感性};
\node[box,below=of p3] (out) {结果表、测线方案\\结论边界};
\draw[arr] (input)--(p1);
\draw[arr] (p1)--(p2);
\draw[arr] (p1)--(p3);
\draw[arr] (p2)--(p4);
\draw[arr] (p3)--(p4);
\draw[arr] (p4)--(check);
\draw[arr] (check)--(out);
\end{tikzpicture}
\caption{四个问题的总体建模与验证流程}
\label{fig:flow}
\end{figure}

\subsection{问题一的分析}
问题一给定的是平面坡面、固定航向和固定测线间距，关键矛盾在于声束两侧面对不同的海底倾斜方向，深水侧和浅水侧的覆盖半宽并不相等。若直接套用平底公式，既无法解释覆盖宽度的非对称性，也无法正确计算相邻条带的交叠。为此，本文先在垂直测线的剖面内建立直角坐标系，通过声束射线与坡面直线的交点求出两侧覆盖半宽，再根据前一条测线浅水边界与当前测线深水边界的相对位置定义有向重叠率。该重叠率允许取负值，因此同一指标既能表示重叠，也能直接识别漏测。

\subsection{问题二的分析}
问题二增加了测线方向角 $\beta$，使海底坡度对水深和覆盖宽度产生两种不同作用：坡度在航向上的投影决定测量船沿程水深变化，坡度在横航向上的投影决定声束扇面内的有效倾角。若不区分这两个分量，就会把沿程变深误当成横向覆盖变宽。本文用单位切向量和单位法向量对深度梯度作正交分解，先计算位置水深，再将横航向坡度代入问题一的剖面公式，由此形成完全可计算的 $D(l,\beta)$、$\gamma(\beta)$ 和 $W(l,\beta)$ 计算链。

\subsection{问题三的分析}
问题三要求在规则矩形海域内满足完整覆盖和 $10\%$--$20\%$ 重叠约束，同时尽量缩短总测线长度。由于等深线沿南北方向，选择南北向直线既能使单条测线沿程水深保持不变，又使每条线取矩形短边长度。本文把候选范围明确限定为“平行等深线的南北向直线族”，在该路线族内以西边界相切确定第一条线，以重叠率下限 $10\%$ 确定后续最大间距，并以首次覆盖东边界为停止条件。由于覆盖半宽与水深均为一次函数，递推式可以显式化，不需要把求解过程隐藏在黑箱优化器中。

\subsection{问题四的分析}
问题四的难点不只是地形复杂，而是同时存在漏测率、过度重叠长度和总测线长度三个目标，且题目没有给出权重。本文采用字典序目标：首先控制漏测，其次减少重叠率超过 $20\%$ 的测线长度，最后比较总长度。真实地形由 $251\times201$ 网格给出，因此需要明确插值公式、梯度公式、路线参数化、覆盖判定和离散精度。候选方案先由保守直线族建立可行基线，再用 $11$ 个横断站位生成地形自适应曲线；原折线存在方向突变，故利用二次 Bézier 圆滑段消除硬拐角，并用曲率半径、交叉数和网格敏感性重新验证。最终方案不是凭单一得分选取，而是在明确的目标层级和失败处理规则下比较得到。

\subsection{模型链归纳}
为使各问所用变量、公式和输出关系一目了然，表~\ref{tab:model-chain} 汇总四个模型的输入、核心状态量、求解方式和直接输出。后续每个小问均按“变量定义—公式建立—算法求解—结果分析—有效性检验”的顺序展开。

\begin{table}[H]
\centering
\caption{四个问题的模型链归纳}
\label{tab:model-chain}
\begin{tabularx}{\textwidth}{C{1.2cm}C{3.0cm}C{3.3cm}C{3.2cm}L}
\toprule
问题 & 主要输入 & 核心状态量与模型 & 求解方法 & 直接输出 \\
\midrule
一 & $D_0,\alpha,\theta,d,x_i$ & 局部水深、两侧覆盖半宽、有向重叠率 & 逐测线显式计算 & 水深、覆盖宽度、重叠率及漏测位置 \\
二 & $D_0,\alpha,\theta,l,\beta$ & 航向深度分量、横航向等效坡度 & $8\times8$ 参数枚举 & 任意方向覆盖宽度矩阵 \\
三 & 矩形边界、$D_0,\alpha,\theta,\eta_0$ & 首线边界方程与显式最大间距递推 & 递推至覆盖东边界 & 测线位置、条数、总长度及路线族内最少性 \\
四 & 深度网格、开角、评价精度 & 双线性地形、局部条带、候选曲线、独立空间指标 & 二分布线、站位插值、Bézier 平滑、离散评估 & 具体测线、漏测率、过度重叠长度、总长度与敏感性 \\
\bottomrule
\end{tabularx}
\end{table}
